\section{结论}
\label{sec:conclusion}

本文研究多智能体量化投资策略，核心发现可以用一句话概括：
在多智能体投资中，复杂的门控策略并不必然优于简单的等权合成。门控的价值等于信号—状态互补性收益减去状态信息的估计成本与交易成本，
只有互补性超过两者之和时，门控才值得做；互补性不足时，
复杂算法只是在为无法利用的状态差异支付估计与交易成本。
这一判断把门控收益的性质界定为信息租金，
并把「何时使用门控」从架构偏好问题转变为成本收益比较问题。

围绕这一判断，本文的理论结果包括：门控占优与否由一个有闭式解的阈值不等式决定；
互补性分解为一阶（收益择时）与二阶（风险配置）两条收敛速率不同的渠道；
制度的马尔可夫持续性通过几何乘数放大门控的信息价值；
容量约束下最优制度数内点存在；满仓长仓约束下以绝对收益为训练目标会使门控退化为
静态等权；横截面宽度不足时互补性在统计上不可识别。
这些结果把「预测组合之谜」从经验规律提升为可计算阈值的判定问题。

在中国 A 股月度面板（868 只证券、247 个月、月均股票池 438 只）上，
实证结果分四层。

训练目标是首个决定成败的因素。未修正时门控退化为静态等权（主动暴露约 $0.0003$，
与等权的夏普差为 $-0.001$、校正后 $p=1.000$），修正后主动暴露升至 0.309。
市场环境决定相对价值的年度分布：年度波动率最低三分位的平均差值为 $+0.006$，
最高三分位为 $+0.189$；横截面离散度呈现同向特征；
市场涨跌方向与差值的关系微弱（秩相关 $-0.049$）。
收益来源以二阶渠道为主：安慰剂检验把主动收益由 $3.77\%$ 压至 $0.004\%$，
说明收益依赖制度信息；Oracle 门控未跑赢静态合成，说明可利用的一阶互补性有限；
两条渠道的状态依赖检验均显著（$p=0.001$ 与 $0.003$）。

边际效应给出三个可操作的参数。交易成本每上升 10 个基点使相对等权的夏普差值
下降约 0.015（$p<0.001$），线性外推的盈亏平衡成本约 92 个基点；
制度状态数呈内点最优，由 2 增至 3 提升 0.079，由 3 增至 4 下降 0.014；
训练窗长度没有单调效应（斜率 $-0.006\,3$/年，$p=0.617$）。
横截面宽度不可调节但决定识别能力，一阶渠道的噪声底线只随时间收敛，
二阶渠道随时间与横截面共同收敛。

稳健性检验界定了结论的适用范围。门控优势在 2020 年后的子样本中由 $+0.180$
收窄至 $+0.049$；四种制度定义下夏普跨度约 0.16，与整体优势同量级；
阈值检验的显著性依赖自助区块长度；十只美国 ETF 的窄横截面上互补性不可识别。
八个可检验含义中五个得到支持，一个不支持，一个为临界结果。

综合判断是：门控组合相对等权取得 $+0.115$ 的夏普差值、3.77\% 的年化主动收益
与 0.80 的主动信息比率，经济量级明显，在 12 个年度观测下统计不显著
（95\% 区间 $[-0.006,0.225]$，Holm 校正后 $p=0.386$）。
多智能体门控运行在可识别性的边界上：其价值真实但脆弱，
更高的成本、更谨慎的推断或更少的样本都会把它推入不显著区间。

把结果放进经济学的语言，本文的证据支持三点判断。
门控收益是状态信息的价格，其大小由互补性与信息成本共同决定：
本文中 $\hat\Gamma$ 与 $\Theta$ 的差距为 $-2.3\times10^{-5}$（一阶）与
$-5.3\times10^{-5}$（二阶），信息的价格与获取成本已经接近。
信息成本的结构解释了渠道差异：一阶矩（预期收益）的估计精度低，
二阶矩（波动与相关结构）可以从公开价格序列低成本估计，
门控收益因此主要来自风险配置，这与理性疏忽与有限套利的预测一致
\cite{sims2003,shleifer1997}。
状态信息的价值随市场分化程度变化：高波动年份与外生冲击年份的门控优势明显更高，
其来源是状态分化幅度而非切换频率，这为「信息在高分化环境中更值钱」提供了直接证据。

对文献的贡献。相对于以架构复杂度为创新点的多智能体与强化学习组合管理研究，本文分离出一个
威胁内部有效性的退化问题：满仓长仓约束下，绝对收益目标中与决策无关的基准项
主导梯度，自适应机制收敛于静态等权。主动暴露可以区分机制奏效与机制退化，
计算成本极低。相对于状态自适应资产配置研究，本文说明状态可识别性本身不充分：
信号有效性不随状态变化时门控价值为零；渠道分解进一步表明，
状态切换的价值可能主要来自风险结构而非收益水平，后者在月度频率上远比前者难以估计。
相对于预测组合文献，本文给出「组合之谜」的临界值形式及其比较静态，
并说明二阶矩比一阶矩更容易估计，以风险配置为主的合成本质上比以收益择时为主的
合成更容易获得稳健收益。相对于人工智能金融治理研究，本文提供三个可复核的诊断量，
把算法审计从原则性要求推进到可执行的判定程序。
在经济学解释上，本文把门控价值界定为信息租金，
并把渠道分解与理性疏忽、有限套利的预测对应起来，
为多智能体策略研究接入信息经济学与行为金融的传统提供了框架。

研究局限与扩展方向。本文结论受月度频率、价量信号与公开数据边界约束，四个扩展方向可以直接检验框架的预测。
其一，提高频率。$\Theta^{(1)}\propto1/T$，日频或周频数据在相同时长内提供更多观测，
门控优势是否随频率上升而变得显著，是对本文框架最直接的检验。
其二，扩充信号集。引入估值、盈利质量或文本情绪信号，
检验专家分歧度上升是否使 $\hat\Gamma$ 越过阈值。
其三，重建数据。使用含完整退市收益与点时成分股的机构数据库重建面板，
消除本文的数据覆盖限制。其四，内生制度识别。用马尔可夫状态切换模型内生估计状态，
检验状态识别误差对阈值判定的净影响。其五，补充微观数据。
本文的证据建立在公开行情与组合层面，无法观测策略的交易对手方与投资者账户，
账户级持仓或订单流数据可以揭示门控收益的微观来源及其对流动性的影响，
这是本文框架之外的扩展方向。

门控组合的研究价值不在于架构的复杂度，而在于它是否在给定的数据与摩擦条件下
满足一个明确的条件。本文的阈值不等式、渠道分解与容量条件使这一判断可以被独立计算与反驳。
报告门控或多智能体组合表现的研究，至少应同时给出三组量：组合相对基准的主动暴露，
用于诊断机制是否退化；互补性统计量与噪声底线的对比，用于诊断机制是否越过阈值；
一阶与二阶渠道的分别检验，用于诊断收益来源。三组量的计算成本都很低。
